\chapter{API Quick Reference}
All routes start with \code{/api}; \code{:name} denotes a path variable. Public requires no token. ADMIN/USER are application roles; Owner permits an administrator or the checked account owner. Signed-in denotes the general authentication requirement. Scopes document source declarations, not executed authorization tests; agent tools require separate execution policy.

\section{Authentication and Directory}
\begingroup\footnotesize\setstretch{1.0}\setlength{\parskip}{0pt}\renewcommand{\arraystretch}{1.15}
\begin{longtable}{>{\raggedright\arraybackslash}p{1.5cm}>{\raggedright\arraybackslash}p{5.9cm}>{\raggedright\arraybackslash}p{1.8cm}>{\raggedright\arraybackslash}p{4.6cm}}
\caption{Authentication and user-directory endpoints.}\\\toprule Method & Route after /api & Scope & Purpose\\\midrule\endfirsthead
\toprule Method & Route after /api & Scope & Purpose\\\midrule\endhead
POST & \code{/auth/register} & Public & Register an active USER; returns 201 and safe account details.\\
POST & \code{/auth/login} & Public & Verify credentials; return JWT and safe account details.\\
GET & \code{/users} & ADMIN & List accounts.\\
POST & \code{/users} & ADMIN & Create an account; returns 201.\\
GET & \code{/users/:id} & ADMIN & Retrieve an account.\\
PUT & \code{/users/:id} & ADMIN & Update account using UserRequestDTO.\\
PATCH & \code{/users/:id/status?active=false} & ADMIN & Set active state; query parameter can be true or false.\\
DELETE & \code{/users/:id} & ADMIN & Delete account; returns 204 on completion.\\\bottomrule
\end{longtable}\endgroup

Registration requires \code{username}, \code{email}, and \code{password}. The DTO limits username to 3--100 characters, email to 150, and password to 8--72. Login requires nonblank username/password and limits them to 100/72 characters. A failed known-account login attempts telemetry recording; an unknown account produces an audit failure without user-linked telemetry.

\section{Telemetry and Anomalies}
\begingroup\footnotesize\setstretch{1.0}\setlength{\parskip}{0pt}\renewcommand{\arraystretch}{1.15}
\begin{longtable}{>{\raggedright\arraybackslash}p{1.5cm}>{\raggedright\arraybackslash}p{5.9cm}>{\raggedright\arraybackslash}p{1.8cm}>{\raggedright\arraybackslash}p{4.6cm}}
\caption{Login-event and anomaly endpoints.}\\\toprule Method & Route after /api & Scope & Purpose\\\midrule\endfirsthead
\toprule Method & Route after /api & Scope & Purpose\\\midrule\endhead
POST & \code{/login-logs} & ADMIN & Persist an explicit event and invoke detection.\\
GET & \code{/login-logs} & ADMIN & List login records.\\
GET & \code{/login-logs/user/:userId} & Owner & Retrieve account-specific logs.\\
GET & \code{/login-logs/failed} & ADMIN & Retrieve failed events.\\
GET & \code{/login-logs/recent?hours=24} & ADMIN & Retrieve recent events; hours defaults to 24.\\
GET & \code{/anomalies} & ADMIN & Paginated list with severity/type filters and sorting.\\
GET & \code{/anomalies/:id} & Owner & Retrieve one anomaly; check its owning account.\\
GET & \code{/anomalies/user/:userId} & Owner & List an account's anomalies.\\
GET & \code{/anomalies/user/:userId/paginated} & Owner & Paginated account anomalies and filters.\\
GET & \code{/anomalies/severity/:severity} & ADMIN & List anomalies by severity.\\
GET & \code{/anomalies/type/:type} & ADMIN & List anomalies by rule type.\\\bottomrule
\end{longtable}\endgroup

Explicit event input requires \code{userId}, \code{ipAddress}, and \code{status}; \code{loginTime}, \code{location}, and \code{device} are optional. Anomaly pagination defaults to \code{page=0}, \code{size=10}, \code{sortBy=detectedAt}, and \code{direction=desc}. Page must be nonnegative, size must be 1--100, and allowed sort fields are \code{id}, \code{detectedAt}, \code{riskScore}, \code{severity}, and \code{anomalyType}.

\section{Risk, Dashboard, and Audit}
\begingroup\footnotesize\setstretch{1.0}\setlength{\parskip}{0pt}\renewcommand{\arraystretch}{1.15}
\begin{longtable}{>{\raggedright\arraybackslash}p{1.5cm}>{\raggedright\arraybackslash}p{5.9cm}>{\raggedright\arraybackslash}p{1.8cm}>{\raggedright\arraybackslash}p{4.6cm}}
\caption{Risk, analytics, dashboard, and audit endpoints.}\\\toprule Method & Route after /api & Scope & Purpose\\\midrule\endfirsthead
\toprule Method & Route after /api & Scope & Purpose\\\midrule\endhead
GET & \code{/risk/user/:userId} & Owner & Account's 24-hour capped risk score and tier.\\
GET & \code{/risk/dashboard} & ADMIN & Overall risk/severity summary.\\
GET & \code{/dashboard/security} & ADMIN & Combined security dashboard.\\
GET & \code{/dashboard/top-risky-users} & ADMIN & Ranked watchlist, up to ten accounts.\\
GET & \code{/dashboard/recent-critical-anomalies} & ADMIN & Recent critical anomalies.\\
GET & \code{/analytics/logins} & Signed-in & Login counts, failure rate, seven-day trend.\\
GET & \code{/audit-logs} & ADMIN & All persisted audit records.\\
GET & \code{/audit-logs/user/:username} & ADMIN & Audit records for a username.\\
GET & \code{/audit-logs/action/:action} & ADMIN & Audit records for an action.\\\bottomrule
\end{longtable}\endgroup

The analytics controller does not declare an ADMIN-only annotation on its login-analytics route; it inherits the authenticated-request requirement from the security chain. API consumers must distinguish this declaration from the administrator-only dashboard routes.

\section{AI Analysis and Agent Execution}
\begingroup\footnotesize\setstretch{1.0}\setlength{\parskip}{0pt}\renewcommand{\arraystretch}{1.15}
\begin{longtable}{>{\raggedright\arraybackslash}p{1.5cm}>{\raggedright\arraybackslash}p{5.9cm}>{\raggedright\arraybackslash}p{1.8cm}>{\raggedright\arraybackslash}p{4.6cm}}
\caption{AI and two-loop agent endpoints.}\\\toprule Method & Route after /api & Scope & Purpose\\\midrule\endfirsthead
\toprule Method & Route after /api & Scope & Purpose\\\midrule\endhead
GET & \code{/ai/status} & ADMIN / USER & Model readiness and mode information.\\
POST & \code{/ai/analyze/anomaly/:anomalyId} & Owner & Dedicated anomaly explanation; alias below.\\
POST & \code{/ai/analyze-anomaly/:anomalyId} & Owner & Alias for anomaly explanation.\\
GET / POST & \code{/ai/analyze/security-summary} & ADMIN & Global security summary; also \code{/ai/summary}.\\
GET & \code{/ai/history} & ADMIN / USER & Paginated history; service scopes ordinary users to their own records.\\
GET & \code{/ai/history/:id} & ADMIN / USER & One history record; service checks cached owner where present.\\
POST & \code{/ai/analyze-user/:userId} & Owner & Dedicated user-risk analysis.\\
POST & \code{/ai/agent-loop} & ADMIN / USER & Execute prompt and queued steering/follow-ups; return trace.\\
GET & \code{/ai/agent-tools} & ADMIN / USER & List registered tools, schemas, roles, and safety metadata.\\
POST & \code{/ai/chat} & ADMIN / USER & Run agent from prompt/context; return simplified chat response.\\\bottomrule
\end{longtable}\endgroup

\subsection{Agent Request Contract}
\begingroup\footnotesize\setstretch{1.0}\renewcommand{\arraystretch}{1.15}
\begin{tabularx}{\textwidth}{p{4.7cm}Y}\toprule Field & Meaning\\\midrule
\code{prompt} & Required nonblank user request.\\
\code{context} & Optional background information.\\
\code{agentRole} & Defaults to DUAL\_COLLABORATIVE; alternatives are THREAT\_INVESTIGATOR and CONTAINMENT\_OPERATOR.\\
\code{steeringMessages} & Optional array of instructions consumed within inner cycles.\\
\code{followUpMessages} & Optional array of continuation requests consumed by the outer loop.\\\bottomrule
\end{tabularx}\endgroup\par
Steering and follow-up arrays are supplied with the request. Agent roles do not replace the authenticated principal's permissions.

\subsection{Response and Error Interpretation}
Successful read/analysis responses normally use HTTP 200. Account/event creation uses 201 where declared; user deletion returns 204. Relevant errors include invalid inputs (400), invalid credentials or tokens (401), denied access/inactive-login rejection (403), missing resources (404), and duplicate registration (409). Error payloads depend on the controller, validation handler, or security filter producing the response.

The agent endpoint returns final text, execution steps, tool names, loop counts, model/mode metadata, and completion time. Availability and fallback indicators must be interpreted alongside actual tool outcomes. A registered operational tool can affect application data even when inference falls back to local planning.
